\documentclass[../main.tex]{subfiles}
\ifSubfilesClassLoaded{\setcounter{chapter}{11}}{}
\begin{document}

\chapter{Pengujian Unit dengan \texttt{unittest} dan \texttt{pytest} (Pengenalan)}

\begin{subcpmk}
  \item Sub-CPMK 7.3: Menulis tes unit sederhana untuk fungsi murni menggunakan \texttt{unittest} atau \texttt{pytest} serta menjalankan tes dari terminal
  \item Sub-CPMK 8.2: Mengonfigurasi \texttt{logging} minimal (level, formatter sederhana) pada skrip kecil agar jejak eksekusi terdokumentasi
\end{subcpmk}

\SesiRPS{13}{7.3, 8.2}{$\sim$170 menit}{CPMK-7, CPMK-8}

\section{Sarana dan prasarana}

Python 3.10+ dengan pustaka standar \texttt{unittest}; \texttt{pytest} opsional (\texttt{pip install pytest} di \texttt{venv} laboratorium jika diizinkan).

\section{Keselamatan kerja dan etika}

Tes harus mencerminkan perilaku yang diharapkan, bukan ``memalsukan'' lolos dengan mengosongkan asersi. Kejujuran akademik berlaku penuh.

\section{Prosedur kerja}

\begin{enumerate}[leftmargin=*]
  \item Pilih 2--3 fungsi murni (tanpa I/O berat) dari latihan sebelumnya atau buat baru.
  \item Tulis berkas \texttt{test\_*.py} dengan kasus normal dan tepi (\textit{edge case}).
  \item Jalankan \texttt{python -m unittest} atau \texttt{pytest -q}; perbaiki kode/fungsi hingga hijau.
  \item Tambahkan \texttt{logging} pada modul yang diuji (mis.\ \texttt{INFO} saat fungsi dipanggil, \texttt{WARNING} pada input tidak wajar); verifikasi keluaran ke konsol atau berkas log sesuai petunjuk dosen.
\end{enumerate}

\section{Tes dengan unittest}

\noindent Kerangka \texttt{unittest} memakai kelas turunan \texttt{TestCase} dan metode yang diawali \texttt{test\_}; asersi seperti \texttt{assertEqual} membandingkan nilai aktual dan harapan. Dokumentasi resmi \href{https://docs.python.org/3/library/unittest.html}{\texttt{unittest}} (lampiran \textbf{[24]}) menyediakan daftar asersi lengkap.

\begin{lstlisting}[caption=unittest dasar]
import unittest

def tambah(a, b):
    return a + b

class TestTambah(unittest.TestCase):
    def test_positif(self):
        self.assertEqual(tambah(2, 3), 5)

    def test_negatif(self):
        self.assertEqual(tambah(-1, 1), 0)

if __name__ == "__main__":
    unittest.main()
\end{lstlisting}

\section{Gaya pytest (opsional)}

\noindent \texttt{pytest} memungkinkan fungsi tes biasa dengan pernyataan \texttt{assert} Python; cocok untuk suite kecil di laboratorium. Panduan resmi \href{https://docs.pytest.org/en/stable/getting-started.html}{Getting Started} (lampiran \textbf{[25]}) menjelaskan penemuan tes otomatis berkas \texttt{test\_*.py}.

\begin{lstlisting}[caption=pytest ringkas]
# test_contoh_pytest.py
def kali(a, b):
    return a * b

def test_kali():
    assert kali(3, 4) == 12
\end{lstlisting}

\noindent Jalankan: \texttt{pytest -q test\_contoh\_pytest.py}.

\section{Contoh tambahan: \texttt{pytest.mark.parametrize}}

\noindent Untuk beberapa kasus uji yang memakai fungsi yang sama, \texttt{parametrize} mengurangi duplikasi. Dokumentasi resmi \href{https://docs.pytest.org/en/stable/how-to/parametrize.html}{pytest parametrizing} (lampiran \textbf{[25]}) menunjukkan pola lanjutan.

\begin{lstlisting}[caption=pytest parametrizasi sederhana]
import pytest

def bagi(a, b):
    return a / b

@pytest.mark.parametrize("a,b,hasil", [(10, 2, 5.0), (9, 3, 3.0)])
def test_bagi(a, b, hasil):
    assert bagi(a, b) == hasil
\end{lstlisting}

\section{\texttt{logging} terstruktur (Sub-CPMK 8.2, pengenalan)}

Modul \texttt{logging} standar menggantikan pola \texttt{print} untuk jejak yang dapat difilter menurut level. Contoh konfigurasi minimal:

\noindent Konfigurasi \texttt{basicConfig} di bawah hanya contoh pengenalan; aplikasi nyata sering memakai logger bernama dan berkas rotasi. Baca \href{https://docs.python.org/3/howto/logging.html}{Logging HOWTO} (lampiran \textbf{[23]}).

\begin{lstlisting}[caption=logging dasar]
import logging

logging.basicConfig(
    level=logging.INFO,
    format="%(levelname)s %(name)s: %(message)s",
)

def bagi(a: float, b: float) -> float:
    logging.info("bagi dipanggil a=%s b=%s", a, b)
    if b == 0:
        logging.warning("pembagi nol")
        raise ValueError("Pembagi tidak boleh nol")
    return a / b
\end{lstlisting}

\noindent Level umum: \texttt{DEBUG}, \texttt{INFO}, \texttt{WARNING}, \texttt{ERROR}, \texttt{CRITICAL}. Untuk uji, Anda dapat menaikkan level log agar konsol tidak berisik, atau mengarahkan ke berkas sementara.

\begin{table}[htbp]
\centering
\small
\begin{tabular}{|l|>{\raggedright\arraybackslash}p{9.8cm}|}
\hline
\textbf{Level} & \textbf{Kegunaan ringkas} \\
\hline
\texttt{DEBUG} & Detail teknis saat pengembangan; biasanya dimatikan di produksi \\
\hline
\texttt{INFO} & Alur normal aplikasi (mis.\ fungsi dipanggil, parameter ringkas) \\
\hline
\texttt{WARNING} & Kondisi tidak wajar namun program tetap jalan \\
\hline
\texttt{ERROR} / \texttt{CRITICAL} & Kegagalan operasi atau kondisi serius \\
\hline
\end{tabular}
\caption{Urutan filter level \texttt{logging} (ringkas). Panduan resmi: \protect\url{https://docs.python.org/3/howto/logging.html} --- \textbf{[23]} pada \texttt{sumber\_materi.md}.}
\label{tab:logging-level}
\end{table}

\section{Studi Kasus dan Contoh Kode}

Berikut adalah contoh-contoh program Python yang mengimplementasikan konsep-konsep pada modul ini.

\subsection{Kode: \texttt{test\_matematika.py}}
\introstudikasuskode
\begin{lstlisting}[language=Python,caption={\texttt{test\_matematika.py}}]
"""Contoh pengujian unit sederhana dengan unittest."""
import unittest


def tambah(a: int, b: int) -> int:
    return a + b

def bagi(a: float, b: float) -> float:
    if b == 0:
        raise ValueError("Pembagi tidak boleh nol")
    return a / b

class TestMatematika(unittest.TestCase):
    def test_tambah(self) -> None:
        self.assertEqual(tambah(2, 3), 5)
        self.assertEqual(tambah(-1, 1), 0)

    def test_bagi(self) -> None:
        self.assertEqual(bagi(10, 2), 5.0)
        with self.assertRaises(ValueError):
            bagi(10, 0)

if __name__ == "__main__":
    unittest.main()
\end{lstlisting}

\section{Referensi sesi}

\begin{itemize}[leftmargin=*]
  \item \textbf{[19]} Real Python --- pengujian dan \texttt{pytest}/\texttt{unittest}.
  \item \textbf{[8]} \textit{The Python Tutorial} --- modul \texttt{unittest}.
  \item \textbf{[13]} \textit{Think Python} --- pengujian dan asersi (bab terkait).
  \item \textbf{[23]} Dokumentasi \texttt{logging} --- konfigurasi dan level.
\end{itemize}

\begin{aktivitas}
  \item Implementasikan fungsi \texttt{palindrome(s)} yang mengembalikan \texttt{True} jika string sama jika dibalik (abaikan spasi opsional sesuai spesifikasi dosen).
  \item Tulis minimal tiga kasus uji (positif, negatif, string kosong) dengan \texttt{unittest}.
  \item Jika prodi mengizinkan \texttt{pytest}, duplikasikan satu kasus dalam sintaks \texttt{assert}.
  \item Tambahkan \texttt{logging} pada modul \texttt{palindrome} (atau berkas utama) dengan minimal satu pesan \texttt{INFO} dan satu \texttt{WARNING}/\texttt{ERROR} pada kasus tidak valid.
  \item Kumpulkan \texttt{NIM\_TIF105\_M13\_test.py} (dan modul yang diuji jika terpisah).
\end{aktivitas}

\begin{latihan}
  \item Apa beda tes formatif saat coding dengan suite \texttt{unittest}?
  \item Mengapa \texttt{if \_\_name\_\_ == "\_\_main\_\_"}\newline sering dipasangkan dengan \texttt{unittest.main()}?
\end{latihan}

\begin{asesmen}
\textbf{Formatif:} suite tes dapat dijalankan; minimal tiga asersi bermakna; satu kasus gagal disengaja kemudian diperbaiki (demonstrasi di laporan singkat); konfigurasi \texttt{logging} terlihat pada jalur eksekusi normal dan kasus error.
\end{asesmen}

\begin{checklist}
  \item Saya dapat menulis kelas turunan \texttt{TestCase}
  \item Saya dapat menjalankan tes dari baris perintah
  \item Saya memahami peran asersi dalam dokumentasi perilaku fungsi
  \item Saya dapat mengonfigurasi \texttt{logging} dasar dan memilih level yang sesuai
\end{checklist}

\begin{rangkuman}
Pengujian unit membantu memastikan fungsi berperilaku sesuai kontrak; \texttt{unittest} built-in; \texttt{pytest} populer untuk sintaks ringkas. Sesi ini memperkenalkan \texttt{logging} sebagai bagian dari CPMK-8 (kualitas operasional kode).

\textbf{Kata kunci:} \code{unittest}, \code{pytest}, asersi, kasus uji, \code{logging}
\end{rangkuman}

\end{document}
